\scanpage{125}
\uline{Dévissage de \(Q_l\)-faisceaux}

\deleted{Un faisceau}

Un \(\unsure{\mathbb Q_l|Q_l}\)-faisceau \deleted{est dit} \inserted{\(F\)} sur \(X\) est dit
\uline{admissible} \inserted{de type \(\le d\)}, s'il satisfait la condition suivante :

\(\forall x\in X\), existe voisinage ouvert \(U\) de \(x\) dans \(\overline{x}\)
et filtration \inserted{finie} de \(F|U\) par des ss-faisceaux
\deleted{\(\unsure{F^i}\)}
tels que les \(\operatorname{gr}^i(F|U)\) soient isomorphes
à un \deleted{``facteur''}
\inserted{\unsure{quotient} \unsure{ess. const.} d'un \unsure{sous-quotient|ss-quotient} \unsure{ess. const.}}
d'un faisceau de la forme
\[
R^ng_*(\deleted{\unsure{Q_{lV}}}\ Q_l),
\]
où \(g:V\mathbin{-}U\) est un morphisme \uline{projectif},
\unsure{et} \uline{lisse}, à fibres connexes \inserted{de dim \(\le d\)}.

\margin{``\unsure{ess. const.}'' signifie : faisceau à \unsure{struc.-disc.|structure discrète}}

[se vérifie fibre par fibre, \(+\ F\) est ``constructible'']

\deleted{\unsure{Cela est} \unsure{que} \(F\) \unsure{est} \illegible{1}}
\inserted{\unsure{en tous cas} \uline{par} \uline{fibre} \uline{géom}}

\uline{Lemme 1}\quad
Les \deleted{\unsure{fi}} \(\unsure{\mathbb Q_l|Q_l}\)-faisceaux admissibles de type \(\le d\)
forment une sous-catégorie abélienne
\deleted{\unsure{exacte} et épaisse (stable par}
\inserted{\uline{épaisse} de la catégorie des}
\inserted{\(Q_l\)-faisceaux constructibles}
\deleted{noyaux, conoyaux, extensions) de la catégorie des}
\deleted{\(\unsure{\mathcal F}\)}
est stable par ss-\deleted{\unsure{fais}}objets et objets quotients.

En effet, si \(F|U\) est muni d'une filtration \(d\)-spéciale,
\inserted{et si \(G|U\) est \unsure{ess. const.},}
alors la filtration induite sur \(G\) ou \((F|U)/G\) est \(d\)-spéciale,
\deleted{car \(H_d\)}
est stable par extensions : trivial.

\uline{Lemme 2}\quad \uline{Images inverses}\hfill trivial.
% archive foliation: 124
